\section{Математические модели в экономике}
\label{sec:models-in-economics}

Экономическая модель выделяет существенные связи между ресурсами,
выпуском, ценами, спросом, предложением и целевыми показателями. В
зависимости от задачи используются балансовые, оптимизационные,
динамические, вероятностные и статистические модели.

\subsection{Межотраслевой баланс Леонтьева}

Пусть $x_i$ --- валовой выпуск отрасли $i$, $y_i$ --- конечный спрос, а
$a_{ij}\ge0$ --- количество продукции отрасли $i$, необходимое для
производства единицы продукции отрасли $j$. Тогда
\[
x_i=\sum_{j=1}^n a_{ij}x_j+y_i.
\]
В матричном виде
\[
\boxed{x=Ax+y},
\qquad
\boxed{(I-A)x=y}.
\]
Если спектральный радиус неотрицательной матрицы прямых затрат удовлетворяет
\[
\boxed{\rho(A)<1},
\]
то матрица $I-A$ обратима и
\[
(I-A)^{-1}=I+A+A^2+\cdots\ge0.
\]
Следовательно, для любого $y\ge0$ существует единственный неотрицательный
вектор выпуска
\[
\boxed{x=(I-A)^{-1}y}.
\]
Это один из основных результатов модели продуктивной экономики Леонтьева.

\subsection{Оптимизация производственного плана}

Пусть предприятие выпускает $n$ видов продукции в количествах
$x_1,\ldots,x_n$. Если $c_j$ --- прибыль с единицы продукции, а
$a_{ij}$ --- расход ресурса $i$ на единицу продукции $j$, то задача
линейного программирования имеет вид
\[
\boxed{
\max_{x\ge0}\;c^Tx
\quad\text{при}\quad
Ax\le b.
}
\]
Здесь $b_i$ --- доступный запас ресурса $i$. Ограничения задают множество
допустимых производственных планов, а целевая функция --- выбранный
экономический критерий.

В линейном программировании при существовании конечного оптимума он
достигается, в частности, в одной из крайних точек допустимого многогранника.
На этом основаны симплексные методы. Двойственные переменные можно
интерпретировать как предельные, или теневые, цены ресурсов.

\subsection{Транспортная задача}

Пусть $a_i$ --- запасы у поставщиков, $b_j$ --- потребности потребителей,
$c_{ij}$ --- стоимость перевозки единицы груза. Необходимо выбрать объёмы
$x_{ij}\ge0$:
\[
\boxed{
\min\sum_{i=1}^m\sum_{j=1}^n c_{ij}x_{ij}
}
\]
при условиях
\[
\sum_jx_{ij}=a_i,
\qquad
\sum_ix_{ij}=b_j.
\]
Сбалансированная задача удовлетворяет
\[
\sum_i a_i=\sum_j b_j.
\]
Если суммы не совпадают, вводят фиктивного поставщика или потребителя.
Транспортная задача является специальным случаем линейного
программирования и имеет более эффективные специализированные алгоритмы.

\subsection{Модель спроса и предложения}

Простейшие линейные зависимости можно записать как
\[
Q_d(p)=a-bp,
\qquad
Q_s(p)=c+dp,
\qquad b,d>0.
\]
Рыночное равновесие определяется условием
\[
Q_d(p^*)=Q_s(p^*),
\]
откуда
\[
\boxed{p^*=\frac{a-c}{b+d}},
\qquad
Q^*=Q_d(p^*)=Q_s(p^*).
\]
Изменение параметров спроса или предложения позволяет исследовать
сравнительную статику равновесия.

\subsection{Максимизация прибыли}

Если предприятие выбирает объём выпуска $q$, цена продукции равна $p$, а
издержки задаются функцией $C(q)$, то
\[
\Pi(q)=pq-C(q).
\]
Для внутреннего оптимума необходимое условие
\[
\boxed{p=C'(q^*)},
\]
то есть цена равна предельным издержкам. Если $C$ строго выпукла,
$C''(q)>0$, то такой стационарный объём является единственным максимумом
прибыли на внутренности допустимого множества.

\subsection{Динамические модели}

Экономические характеристики могут изменяться во времени. Например,
простейшая модель накопления капитала имеет вид
\[
\dot K=sF(K)-\delta K,
\]
где $F(K)$ --- производственная функция, $s$ --- норма сбережений,
$\delta$ --- коэффициент выбытия. Стационарные состояния удовлетворяют
\[
sF(K^*)=\delta K^*.
\]
Такие модели исследуют методами теории динамических систем: находят
равновесия, изучают их устойчивость и влияние параметров.


\subsection{Двойственная задача и экономический смысл цен ресурсов}

Для задачи
\[
    \max\;c^Tx,
    \qquad Ax\le b,
    \qquad x\ge0
\]
двойственная задача имеет вид
\[
    \boxed{
    \min\;b^T\lambda,
    \qquad A^T\lambda\ge c,
    \qquad \lambda\ge0.
    }
\]
Для любых допустимых $x$ и $\lambda$ выполняется слабая двойственность
\[
    c^Tx\le b^T\lambda.
\]
При стандартных условиях разрешимости линейного программирования оптимальные
значения прямой и двойственной задач совпадают. Условия дополняющей нежёсткости
имеют вид
\[
    \lambda_i\bigl(b_i-(Ax)_i\bigr)=0,
\]
\[
    x_j\bigl((A^T\lambda)_j-c_j\bigr)=0.
\]
Компонента $\lambda_i$ интерпретируется как предельная ценность дополнительной
единицы $i$-го ресурса: если ограничение не активно, его теневая цена равна нулю.

\subsection{Интерпретация обратной матрицы Леонтьева}

Разложение
\[
    (I-A)^{-1}=I+A+A^2+\cdots
\]
имеет прямой экономический смысл. Матрица $A$ описывает непосредственные затраты,
$A^2$ --- затраты ресурсов второго уровня через цепочки поставок, $A^3$ ---
третьего уровня и т.д. Поэтому элемент обратной матрицы Леонтьева характеризует
полную прямую и косвенную потребность отраслей, необходимую для обеспечения
единицы конечного спроса.

При малом изменении конечного спроса
\[
    \Delta x=(I-A)^{-1}\Delta y,
\]
что позволяет проводить анализ чувствительности балансовой экономики.

\subsection{Двойственность транспортной задачи}

Для сбалансированной транспортной задачи двойственные переменные можно связать
с потенциалами поставщиков $u_i$ и потребителей $v_j$. Двойственные ограничения
имеют вид
\[
    u_i+v_j\le c_{ij}.
\]
Если $x_{ij}>0$ в оптимальном плане, то по дополняющей нежёсткости
\[
    u_i+v_j=c_{ij}.
\]
Эта структура лежит в основе метода потенциалов и показывает, почему специальная
структура транспортной задачи позволяет решать её эффективнее общего ЛП.

\subsection{Производственная функция Кобба--Дугласа}

Распространённая нелинейная модель выпуска имеет вид
\[
    \boxed{Y=A K^\alpha L^\beta,}
\]
где $K$ --- капитал, $L$ --- труд, $A$ --- технологический коэффициент.
Эластичности выпуска равны
\[
    \frac{\partial\ln Y}{\partial\ln K}=\alpha,
    \qquad
    \frac{\partial\ln Y}{\partial\ln L}=\beta.
\]
Сумма $\alpha+\beta$ характеризует отдачу от масштаба: при значении $1$ она
постоянная, при значении больше $1$ --- возрастающая, меньше $1$ --- убывающая.
Параметры такой модели могут оцениваться по данным после логарифмирования или
нелинейным МНК.

\subsection{Локальная устойчивость динамической экономической модели}

Для
\[
    \dot K=g(K)=sF(K)-\delta K
\]
стационарное состояние $K^*$ удовлетворяет $g(K^*)=0$. Линеаризация
\[
    \dot{\kappa}\approx g'(K^*)\kappa,
    \qquad \kappa=K-K^*,
\]
показывает, что равновесие локально асимптотически устойчиво при
\[
    \boxed{g'(K^*)<0.}
\]
Таким образом, динамическая экономическая модель исследуется теми же средствами,
что и другие нелинейные динамические системы: равновесия, линеаризация,
устойчивость и бифуркации.

\subsection{Параметры, данные и проверка модели}

Экономические коэффициенты обычно не известны заранее точно и оцениваются по
наблюдаемым данным. Поэтому кроме решения уравнений необходимо проверять
идентифицируемость параметров, чувствительность прогноза и устойчивость вывода к
изменению периода наблюдений. Хорошее описание исторических данных не гарантирует
прогноз вне использованного диапазона: структурные изменения экономики могут
сделать прежние параметры неприменимыми.

\subsection{Что важно сказать на экзамене}

Лучше показать несколько принципиально разных классов: \textbf{балансовая
модель Леонтьева}, \textbf{линейное программирование и транспортная задача},
\textbf{равновесие спроса и предложения} и простую динамическую модель.
Для Леонтьева полезно записать условие $\rho(A)<1$ и формулу
$(I-A)^{-1}=I+A+A^2+\cdots$. В оптимизационных моделях обязательно назвать
переменные, ограничения и целевую функцию.

\newpage
